\section{既有实验结果的综合}
\label{sec:prior-findings}

条件比较不仅保存数值，也能识别不同研究中反复出现的规律。综合目标相与相关化合物的已发表实验，可归纳出三项可直接指导下一轮实验的认识。

\subsection{相形成差异的主导因素：组成坐标而非单一最高温度}

Ba--Y--Si--O体系助熔剂法生长的151组实验表明，Y$_2$O$_3$\slash SiO$_2$比和MoO$_3$用量主导硅酸根的连接方式与晶体类型。降低Y$_2$O$_3$\slash SiO$_2$比或在一定范围内增加MoO$_3$用量，结构均趋向更高的SiO$_4$连接度；MoO$_3$过量则不利于晶体生长，较低的冷却速率通常改善晶体形貌和产率\cite{wierzbickawieczorek2017high}。因此，目标相筛查应将营养物组成、助熔剂比例和冷却速率视为独立变量，而不能仅复制某一最高温度。早期伴生晶体在1150\,$^\circ$C保温并缓冷后获得，但该结果来自含K$_2$CO$_3$\slash MoO$_3$的整批助熔体系，不能换算为目标相的化学计量投料\cite{kolitsch2009crystal}。

固相研究指向同一结论。$\mathrm{Ba_xY_{26}Si_{16}O_{71+x}}$陶瓷在$9\le x\le 14$组成系列中，主相与第二相随名义Ba含量变化；玛瑙研磨带入的SiO$_2$会干扰本征相区的判断\cite{yamane2024microstructure}。$\mathrm{Ba_5Y_{13}[SiO_4]_8O_{8.5}}$是在跨越相图的粉体组成筛查中定位，并经反复退火、研磨以及多种衍射与成分分析方法确认\cite{gulay2024navigation}。因此，目标相目前缺少的不是又一个孤立温度，而是Ba--Y--Zn--Si--O局部相图，以及与每个组成坐标对应的定量相分析。

\subsection{单晶证据与块体相纯窗口的区别}

目标相的直接报道以单晶衍射确认了结构，但现有资料未同时给出可独立复现的投料、热历史和块体相分数\cite{ababaikeri2024ba5y12zn}。无Zn系列的陶瓷研究则表明，即使已形成主相，块体样品仍可能含有第二相和处理介质带入的污染\cite{yamane2024microstructure,gulay2024navigation}。因此，复现实验至少需要两类互补表征：以单晶或先进衍射方法确认结构模型，以定量PXRD\slash Rietveld测定整批样品中的目标相比例。二者缺一，“获得目标晶体”即不能等同于“建立相纯制备方法”。

\subsection{局部相图：连接复现与发现的核心实验}

$\mathrm{Ba_5Y_{13}[SiO_4]_8O_{8.5}}$的发现表明，即使BaO--Y$_2$O$_3$--SiO$_2$相区已知多种化合物，新相仍可能位于竞争相或玻璃区的边界。研究者依据总体组成、已知相比例和成分分析逐轮更新下一批投料，才逐步缩小目标区域\cite{gulay2024navigation}。含Zn体系更需要这一策略，因为Zn的引入会同时改变阳离子数、氧计量以及可能的液相行为。相图实验可区分三种情形：窄组成的线化合物、无Zn系列中的有限固溶范围，以及仅沿特定液相路径形成的亚稳晶相。与反复尝试单一温度相比，这三种结果均能提供更多的结构信息。
